%%% آپدیت شده در شهریور 1402

%%% مواردی که نسبت به قبل بروزرسانی شده است به شرح زیر می باشد
%%% 1- بولد شدن سر فصل ها
%%% 2- اضافه شدن قابلیت
%%%    \begin{subfigure}
%%% 3-   اصلاح ترتیب شماره گذاری مراجع به ترتیب اولویت ارجاع در متن 
%%% 4- مرتب شدن فایل‌ها (قرار گرفتن عکس ها و فونت ها در پوشه ای جداگانه)
%%% 5- اضافه شدن پکیج فرض
%%% 6- اضافه شدن فونت یاس برای فارسی شدن اعداد
%%% 7- اضافه شدن پکیج تنظیم سایز جدول


%%%  کلاس AUTthesis، نسخه آبان 1397
%%%   دانشگاه صنعتی امیرکبیر                 http://www.aut.ac.ir
%%%  تالار گفتگوی پارسی‌لاتک،       http://forum.parsilatex.com
%%%   آپدیت شده در آبان 95
%%%   پشتیبانی و راهنمایی          badali_farhad@yahoo.com
%%%
%%%   بازبینی و اصلاح شده در آبان ماه 1397
%%%  Tested via TeXstudio in TeXlive 2014-2018.
%%%

%-----------------------------------------------------------------------------------------------------
%        روش اجرا.: 2 بار F1 ، 2 بار  F11(به منظور تولید مراجع) ، دوبار Ctrl+Alt+I (به منظور تولید نمایه) و دو بار F1 -------> مشاهده Pdf
%%%%%%%%%%%%%%%%%%%%%%%%%%%%%%%%%%%%%%%%%%%%%%%%%%%%%%
%   TeXstudio as your IDE
%%  برای compile در TeXstudio تنها کافی است منوی Options->Configure TeXstudio را زده و در پنجره Configure TeXstudio در بخش Build گزینه Default Compiler را به XeLaTeX تغییر دهید. سند شما به راحتی compile خواهد شد.
%   F1 & F5 : Build & view
%   F6      : Compile
%   F7      : View
%   --------------
%%%%%%%%%%%%%%%%%%%%%%%%%%%%%%%%%%%%%%%%%%%%%%%%%%%%%%
%        اگر قصد نوشتن رساله دکتری را دارید، در خط زیر به جای msc،
%      کلمه phd را قرار دهید. کلیه تنظیمات لازم، به طور خودکار، اعمال می‌شود.
%%% !TEX TS-program = XeLaTeX
\documentclass[oneside,msc,12pt]{AUTthesis}
%       فایل commands.tex را حتماً به دقت مطالعه کنید؛ چون دستورات مربوط به فراخوانی بسته زی‌پرشین 
%       و دیگر بسته‌ها و ... در این فایل قرار دارد و بهتر است که با نحوه استفاده از آنها آشنا شوید. توجه شود برای نسخه نهایی پایان‌نامه حتماً hyperref را 
%        غیرفعال کنید.


% در این فایل، دستورها و تنظیمات مورد نیاز، آورده شده است.
%-------------------------------------------------------------------------------------------------------------------
% در ورژن جدید زی‌پرشین برای تایپ متن‌های ریاضی، این سه بسته، حتماً باید فراخوانی شود.
\usepackage{amsthm,amssymb,amsmath,amsfonts}
% بسته‌ای برای تنطیم حاشیه‌های بالا، پایین، چپ و راست صفحه
\usepackage[top=30mm, bottom=30mm, left=25mm, right=30mm]{geometry}
% بسته‌‌ای برای ظاهر شدن شکل‌ها و تصاویر متن
\usepackage{graphicx}
\usepackage{color}
%بسته‌ای برای تنظیم فاصله عمودی خط‌های متن
\usepackage{setspace}
\usepackage{titletoc}
\usepackage{tocloft}
%با فعال کردن بسته زیر فوت‌نوت‌ها در هر صفحه ریست می‌شوند. حالت پیش‌فرض آن ریست شدن در هر فصل می‌باشد.
%\usepackage[perpage]{footmisc}
\usepackage{enumitem}
\usepackage{multirow,adjustbox}
%\usepackage{titlesec}
% بسته‌ و دستوراتی برای ایجاد لینک‌های رنگی با امکان جهش
\usepackage[pagebackref=false,colorlinks,linkcolor=blue,citecolor=red]{hyperref}
\usepackage[nameinlink]{cleveref}%capitalize,,noabbrev
 \AtBeginDocument{%
    \crefname{equation}{برابری}{equations}%
    \crefname{chapter}{فصل}{chapters}%
    \crefname{section}{بخش}{sections}%
    \crefname{appendix}{پیوست}{appendices}%
    \crefname{enumi}{مورد}{items}%
    \crefname{footnote}{زیرنویس}{footnotes}%
    \crefname{figure}{شکل}{figures}%
    \crefname{table}{جدول}{tables}%
    \crefname{theorem}{قضیه}{theorems}%
    \crefname{lemma}{لم}{lemmas}%
    \crefname{corollary}{نتیجه}{corollaries}%
    \crefname{proposition}{گزاره}{propositions}%
    \crefname{definition}{تعریف}{definitions}%
    \crefname{result}{نتیجه}{results}%
    \crefname{example}{مثال}{examples}%
    \crefname{remark}{نکته}{remarks}%
    \crefname{note}{یادداشت}{notes}%
    \crefname{asum}{فرض}{Assumption}
    % دستوری برای تغییر نام کلمه «کتاب‌نامه» به «منابع و مراجع«
    \renewcommand{\bibname}{منابع و مراجع}
    
    \renewcommand{\labelitemi}{$\bullet$}
    % دستوری برای تعیین علامت سطح دوم itemize
    \renewcommand{\labelitemii}{$\circ$}
    % دستوری برای تعیین علامت سطح سوم itemize
    \renewcommand{\labelitemiii}{$-$}
    % برای سطح چهارم
    \renewcommand{\labelitemiv}{$*$}
    % دستوری برای تغییر نام کلمه «اثبات» به «برهان»
    \renewcommand\proofname{\textbf{برهان}}
    \renewcommand{\listfigurename}{فهرست اشکال}
    \renewcommand{\listtablename}{فهرست جداول}
}

\usepackage{subcaption}
% چنانچه قصد پرینت گرفتن نوشته خود را دارید، خط بالا را غیرفعال و  از دستور زیر استفاده کنید چون در صورت استفاده از دستور زیر‌‌، 
% لینک‌ها به رنگ سیاه ظاهر خواهند شد که برای پرینت گرفتن، مناسب‌تر است
%\usepackage[pagebackref=false]{hyperref}
% بسته‌ لازم برای تنظیم سربرگ‌ها
\usepackage{fancyhdr}
% بسته‌ای برای ظاهر شدن «مراجع»  در فهرست مطالب
\usepackage[nottoc]{tocbibind}
% دستورات مربوط به ایجاد نمایه
\usepackage{makeidx,multicol}
\setlength{\columnsep}{1.5cm}

%%%%%%%%%%%%%%%%%%%%%%%%%%
\usepackage{verbatim}
\makeindex
\usepackage{sectsty}
% فراخوانی بسته زی‌پرشین و تعریف قلم فارسی و انگلیسی
\usepackage{xepersian}%[extrafootnotefeatures]
\SepMark{-}
%حتماً از تک لایو 2014 استفاده کنید.
\settextfont[Scale=1.2,Path=Fonts/,BoldFont=B Nazanin Bold.ttf]{B Nazanin.ttf}
\setlatintextfont[Path=Fonts/,BoldFont=timesbd]{times}

%%%%%%%%%%%%%%%%%%%%%%%%%%
% چنانچه می‌خواهید اعداد در فرمول‌ها، انگلیسی باشد، خط زیر را غیرفعال کنید.
%
\setdigitfont[Scale=1.1,Path=Fonts/,BoldFont=Yas Bd.ttf]{Yas.ttf}%%Yas
%%%%%%%%%%%%%%%%%%%%%%%%%%
% تعریف قلم‌های فارسی اضافی برای استفاده در بعضی از قسمت‌های متن
\defpersianfont\nastaliq[Scale=2,Path=Fonts/]{IranNastaliq.ttf}
\defpersianfont\chapternumber[Scale=3,Path=Fonts/,BoldFont=B Nazanin Bold.ttf]{B Nazanin.ttf}
%\chapterfont{\centering}%
%%%%%%%%%%%%%%%%%%%%%%%%%%





% Headings for every page of ToC, LoF and Lot
\setlength{\cftbeforetoctitleskip}{-1.2em}
\setlength{\cftbeforelottitleskip}{-1.2em}
\setlength{\cftbeforeloftitleskip}{-1.2em}
\setlength{\cftaftertoctitleskip}{-1em}
\setlength{\cftafterlottitleskip}{-1em}
\setlength{\cftafterloftitleskip}{-1em}
%%\makeatletter
%%%%\renewcommand{\l@chapter}{\@dottedtocline{1}{1em\bfseries}{1em}}
%%%%\renewcommand{\l@section}{\@dottedtocline{2}{2em}{2em}}
%%%%\renewcommand{\l@subsection}{\@dottedtocline{3}{3em}{3em}}
%%%%\renewcommand{\l@subsubsection}{\@dottedtocline{4}{4em}{4em}}
%%%%\makeatother


\newcommand\tocheading{\par عنوان\hfill صفحه \par}
\newcommand\lofheading{\hspace*{.5cm}\figurename\hfill صفحه \par}
\newcommand\lotheading{\hspace*{.5cm}\tablename\hfill صفحه \par}

\renewcommand{\cftchapleader}{\cftdotfill{\cftdotsep}}
\renewcommand{\cfttoctitlefont}{\hspace*{\fill}\LARGE\bfseries}%\Large
\renewcommand{\cftaftertoctitle}{\hspace*{\fill}}
\renewcommand{\cftlottitlefont}{\hspace*{\fill}\LARGE\bfseries}%\Large
\renewcommand{\cftafterlottitle}{\hspace*{\fill}}
\renewcommand{\cftloftitlefont}{\hspace*{\fill}\LARGE\bfseries}
\renewcommand{\cftafterloftitle}{\hspace*{\fill}}

%%%%%%%%%%%%%%%%%%%%%%%%%%
% تعریف و نحوه ظاهر شدن عنوان قضیه‌ها، تعریف‌ها، مثال‌ها و ...
%برای شماره گذاری سه تایی قضیه ها
\theoremstyle{definition}
\newtheorem{definition}{تعریف}[section]
\newtheorem{remark}[definition]{نکته}
\newtheorem{note}[definition]{یادداشت}
\newtheorem{example}[definition]{نمونه}
\newtheorem{question}[definition]{سوال}
\newtheorem{remember}[definition]{یاداوری}
\theoremstyle{theorem}
\newtheorem{theorem}[definition]{قضیه}
\newtheorem{lemma}[definition]{لم}
\newtheorem{proposition}[definition]{گزاره}
\newtheorem{corollary}[definition]{نتیجه}
\newtheorem{asum}[definition]{فرض}
%%%%%%%%%%%%%%%%%%%%%%%%
%%%%%%%%%%%%%%%%%%%
%%% برای شماره گذاری چهارتایی قضیه ها و ...
%%\newtheorem{definition1}[subsubsection]{تعریف}
%%\newtheorem{theorem1}[subsubsection]{قضیه}
%%\newtheorem{lemma1}[subsubsection]{لم}
%%\newtheorem{proposition1}[subsubsection]{گزاره}
%%\newtheorem{corollary1}[subsubsection]{نتیجه}
%%\newtheorem{remark1}[subsubsection]{نکته}
%%\newtheorem{example1}[subsubsection]{مثال}
%%\newtheorem{question1}[subsubsection]{سوال}

%%%%%%%%%%%%%%%%%%%%%%%%%%%%

% دستورهایی برای سفارشی کردن صفحات اول فصل‌ها
\makeatletter
\newcommand\mycustomraggedright{%
 \if@RTL\raggedleft%
 \else\raggedright%
 \fi}
\def\@makechapterhead#1{%
\thispagestyle{style1}
\vspace*{20\p@}%
{\parindent \z@ \mycustomraggedright
\ifnum \c@secnumdepth >\m@ne
\if@mainmatter

\bfseries{\Huge \@chapapp}\small\space {\chapternumber\thechapter}
\par\nobreak
\vskip 0\p@
\fi
\fi
\interlinepenalty\@M 
\Huge \bfseries #1\par\nobreak
\vskip 120\p@

}

%\thispagestyle{empty}
\newpage}
\bidi@patchcmd{\@makechapterhead}{\thechapter}{\tartibi{chapter}}{}{}
\bidi@patchcmd{\chaptermark}{\thechapter}{\tartibi{chapter}}{}{}
\makeatother

\pagestyle{fancy}
\renewcommand{\chaptermark}[1]{\markboth{\chaptername~\tartibi{chapter}: #1}{}}

\fancypagestyle{style1}{
\fancyhf{} 
\fancyfoot[c]{\thepage}
\fancyhead[R]{\leftmark}%
\renewcommand{\headrulewidth}{1.2pt}
}


\fancypagestyle{style2}{
\fancyhf{}
\fancyhead[R]{چکیده}
\fancyfoot[C]{\thepage{}}
\renewcommand{\headrulewidth}{1.2pt}
}

\fancypagestyle{style3}{%
  \fancyhf{}%
  \fancyhead[R]{فهرست نمادها}
  \fancyfoot[C]{\thepage}%
  \renewcommand{\headrulewidth}{1.2pt}%
}

\fancypagestyle{style4}{%
  \fancyhf{}%
  \fancyhead[R]{فهرست جداول}
  \fancyfoot[C]{\thepage}%
  \renewcommand{\headrulewidth}{1.2pt}%
}

\fancypagestyle{style5}{%
  \fancyhf{}%
  \fancyhead[R]{فهرست اشکال}
  \fancyfoot[C]{\thepage}%
  \renewcommand{\headrulewidth}{1.2pt}%
}

\fancypagestyle{style6}{%
  \fancyhf{}%
  \fancyhead[R]{فهرست مطالب}
  \fancyfoot[C]{\thepage}%
  \renewcommand{\headrulewidth}{1.2pt}%
}

\fancypagestyle{style7}{%
  \fancyhf{}%
  \fancyhead[R]{نمایه}
  \fancyfoot[C]{\thepage}%
  \renewcommand{\headrulewidth}{1.2pt}%
}

\fancypagestyle{style8}{%
  \fancyhf{}%
  \fancyhead[R]{منابع و مراجع}
  \fancyfoot[C]{\thepage}%
  \renewcommand{\headrulewidth}{1.2pt}%
}
\fancypagestyle{style9}{%
  \fancyhf{}%
  \fancyhead[R]{واژه‌نامه‌ی فارسی به انگلیسی}
  \fancyfoot[C]{\thepage}%
  \renewcommand{\headrulewidth}{1.2pt}%
}
%


%دستور حذف نام لیست تصاویر و لیست جداول از فهرست مطالب
\newcommand*{\BeginNoToc}{%
  \addtocontents{toc}{%
    \edef\protect\SavedTocDepth{\protect\the\protect\value{tocdepth}}%
  }%
  \addtocontents{toc}{%
    \protect\setcounter{tocdepth}{-10}%
  }%
}
\newcommand*{\EndNoToc}{%
  \addtocontents{toc}{%
    \protect\setcounter{tocdepth}{\protect\SavedTocDepth}%
  }%
}
\newcounter{savepage}

%\renewcommand\cftsecleader{\cftdotfill{\cftdotsep}}
%%%%%%%%%%%%%%%%%%%%%%%%%%%%%
%%%%%%%%%%%%%%%%%%%%%%%%%%%%




\begin{document}
\baselineskip=.75cm
\linespread{1.75}
%% -!TEX root = AUTthesis.tex
% در این فایل، عنوان پایان‌نامه، مشخصات خود، متن تقدیمی‌، ستایش، سپاس‌گزاری و چکیده پایان‌نامه را به فارسی، وارد کنید.
% توجه داشته باشید که جدول حاوی مشخصات پروژه/پایان‌نامه/رساله و همچنین، مشخصات داخل آن، به طور خودکار، درج می‌شود.
%%%%%%%%%%%%%%%%%%%%%%%%%%%%%%%%%%%%
% دانشکده، آموزشکده و یا پژوهشکده  خود را وارد کنید
\faculty{دانشکده ریاضی و علوم کامپیوتر}
% گرایش و گروه آموزشی خود را وارد کنید
\department{رشته علوم کامپیوتر}
% عنوان پایان‌نامه را وارد کنید
\fatitle{
طبقه بندی و خوشه بندی داده
	\\[.75 cm]
	گزارش تمرین سری پنجم}
% نام استاد(ان) راهنما را وارد کنید
\firstsupervisor{دکتر مهدی قطعی}
%\secondsupervisor{استاد راهنمای دوم}
% نام استاد(دان) مشاور را وارد کنید. چنانچه استاد مشاور ندارید، دستور پایین را غیرفعال کنید.
%\firstadvisor{دکتر مهدی قطعی}
%\secondadvisor{استاد مشاور دوم}
% نام نویسنده را وارد کنید
\name{محمدرضا}
% نام خانوادگی نویسنده را وارد کنید
\surname{امانی}
%%%%%%%%%%%%%%%%%%%%%%%%%%%%%%%%%%
% \thesisdate{\date}

% چکیده پایان‌نامه را وارد کنید
\fa-abstract{
مجموعه داده شامل بیش از 284807 داده از داده‌‌های بانک‌های اروپا شامل بیش از ۲۵ کلاس و همچنین زمان و مقدار آن‌هاست که در اینجا با 1 و 0 شماره گذاری شده اند.
ما در اینجا قصد داریم ابتدا با پیش پردازش و سپس با استفاده از الگوریتم \lr{kmeans} برای خوشه بندی و همچنین الگوریتم های \lr{SVC, LogisticRegression, KNeighborsClassifier} 
به پاسخی برای دسته بندی و همچنین خوشه بندی مجموعه داده برسیم.	
}


% کلمات کلیدی پایان‌نامه را وارد کنید
\keywords{
	\lr{classification}
	، هوش مصنوعی،
	\lr{clustering}
}



% \AUTtitle
%%%%%%%%%%%%%%%%%%%%%%%%%%%%%%%%%%
\vspace*{7cm}
\thispagestyle{empty}
\begin{center}
	\includegraphics[height=5cm,width=12cm]{Images/besm.jpg}
\end{center}
% تاییدیه دفاع
%\newpage
\thispagestyle{empty}
%\fontsize{18pt}{19pt}\selectfont

\section*{صفحه فرم ارزیابی و تصویب پایان نامه- فرم تأیید اعضاء كميته دفاع}

\fontsize{12pt}{14pt}\selectfont
%\renewcommand{\baselinestretch}{1.5}
\vspace*{1cm}
   در این صفحه فرم دفاع یا تایید و تصویب پایان نامه موسوم به فرم کمیته دفاع- موجود در پرونده آموزشی- را قرار دهید.
\vspace*{1cm}


\subsection*{نکات مهم:}
 
\begin{itemize}
\item
	نگارش پایان نامه/رساله باید به
	{\color{red}
		زبان فارسی
	}
	و بر اساس آخرین نسخه دستورالعمل و راهنمای تدوین پایان نامه های دانشگاه صنعتی امیرکبیر باشد.(دستورالعمل و راهنمای حاضر)
\item رنگ جلد پایان نامه/رساله چاپي كارشناسي، كارشناسي ارشد و دكترا  بايد به ترتيب مشكي، طوسي و سفيد رنگ باشد.  
\item چاپ و صحافی پایان نامه/رساله بصورت
{\color{red}
	پشت و رو(دورو)
}
بلامانع است و انجام آن توصيه مي شود. 
\end{itemize}
%%%%%%%%%%%%%%%%%%%%%%%%%%%%%%%%%%%%%%%%%%%%%%%%%%%%%%%%%%%%%%%%%%%%%%%%%%%%%%%%%%%%%%%%%%%%%%%%%%
%%%%%%%%%%%%%%%%%%%%%%%%%%%%%%%%%%%%%%%%%%%%%%%%%%%%%%%%%%%%%%%%%%%%%%%%%%%%%%%%%%%%%%%%%%%%%%%%%%
\newpage
\thispagestyle{empty}
\begin{picture}(50,50)
  \put(17,0){\includegraphics[scale=1.1]{fa-logo.png}}
  \put(4.5,-13){\footnotesize{دانشگاه صنعتی امیرکبیر}}
  \put(10.5,-27){\footnotesize{(پلی‌تکنیک تهران)}}
  \put(170,30){\bf{به نام خدا}}
  \put(140,-5){\Large\bf{تعهدنامه اصالت اثر}}
  \put(310,0){تاریخ: \datethesis}
\end{picture}

\vspace*{2.5cm}

اينجانب {\bf{\fname\lname}} متعهد می‌شوم که مطالب مندرج در این پایان‌نامه حاصل کار پژوهشی اینجانب تحت نظارت و راهنمایی اساتید دانشگاه صنعتی امیرکبیر بوده و به دستاوردهای دیگران که در این پژوهش از آنها استفاده شده است مطابق مقررات و روال متعارف ارجاع و در فهرست منابع و مآخذ ذکر گردیده است. این پایان‌نامه قبلاً برای احراز هیچ مدرک هم‌سطح یا بالاتر ارائه نگردیده است.

در صورت اثبات تخلف در هر زمان، مدرک تحصیلی صادر شده توسط دانشگاه از درجه اعتبار ساقط بوده و دانشگاه حق پیگیری قانونی خواهد داشت.


کلیه نتایج و حقوق حاصل از این پایان‌نامه متعلق به دانشگاه صنعتی امیرکبیر می‌باشد. هرگونه استفاده از نتایج علمی و عملی، واگذاری اطلاعات به دیگران یا چاپ و تکثیر، نسخه‌برداری، ترجمه و اقتباس از این پایان نامه بدون موافقت کتبی دانشگاه صنعتی امیرکبیر ممنوع است. 
نقل مطالب با ذکر مآخذ بلامانع است.\\
\vspace{2.5cm}


{\centerline {\bf{\fname\lname}}}
\vspace*{.2cm}
{\centerline{امضا}}
%%%%%%%%%%%%%%%%%%%%%%%%%%%%%%%%%
% چنانچه مایل به چاپ صفحات «تقدیم»، «نیایش» و «سپاس‌گزاری» در خروجی نیستید، خط‌های زیر را با گذاشتن ٪  در ابتدای آنها غیرفعال کنید.
% پایان‌نامه خود را تقدیم کنید
% نیایش خود را در فایل زیر بنویسید.
%\begin{acknowledgementpage}

\vspace{1.5cm}

{\nastaliq
{
}}\end{acknowledgementpage}
\newpage
% سپاسگزاری را در فایل زیر بنویسید.
%%%%%%%%%%%%%%%%%%%%%%%%%%%%%%%%%%%%%%%%%
%%%%%%%%%%%%%%%%%%%%%%%%%%%%%%%%%کدهای زیر را تغییر ندهید.
\newpage\clearpage

\pagestyle{style2}

\vspace*{-1cm}
\section*{\centering چکیده}
%\addcontentsline{toc}{chapter}{چکیده}
\vspace*{.5cm}
\ffa-abstract
\vspace*{2cm}


{\noindent\large\textbf{واژه‌های کلیدی:}}\par
\vspace*{.5cm}
\fkeywords
% دستور زیر برای شماره گذاری صفحات قبل از فصل اول با حروف ابجد است.
\pagenumbering{alph}
%-----------------------------------------------------------------------------
% فایل زیر دستورات مربوط به نمایش صفحات فهرست مطالب- فهرست اشکال و جداول است.
\newpage

%{\pagestyle{style2}
%\tableofcontents}\newpage

%\listoffigures
%\cleardoublepage
 \pagestyle{style4}
\tableofcontents
 \pagestyle{style4}
\cleardoublepage
%اگر لیست تصاویر و لیست جداول ندارید ، کدهای زیر را با گذاشتن % در ابتدای آنها، غیرفعال کنید.
%\BeginNoToc
%%============
\addtocontents{lof}{\lofheading}% add heading to the first page in LoF
\pagestyle{style5}
\listoffigures
\thispagestyle{style5}
\cleardoublepage
%%============
%% \addtocontents{lot}{\lotheading}% add heading to the first page in LoT
%% \thispagestyle{style4}
%% \listoftables
%% \thispagestyle{style4}
%
%%============
%\cleardoublepage
%
%\cleardoublepage
%\setcounter{savepage}{\arabic{page}}
%\mainmatter
%\addtocontents{toc}{\tocheading}% add heading to the first page in ToC, after frontmatter entries
%\EndNoToc

% در صورت تمایل می‌توانید با فعال کردن دستور بالا، لیست تصاویر را به  پایان‌نامه خود اضافه کنید.
%-------------------------------------------------------------------------symbols(فهرست نمادها)
% وجود لیست نمادها الزامیست.(لطفاً نمادهای خود را جایگذین نمادهای پیش‌فرض کنید.)
% %%%%%%%%%%%%%

{\centering\LARGE\textbf{فهرست نمادها}\par}%

\pagenumbering{alph}
\setcounter{page}{\thesavepage}
%\setcounter{page}{6}
\vspace*{1cm}

\pagestyle{style3}
%\thispagestyle{empty}
%\addcontentsline{toc}{chapter}{فهرست نمادها}
\symb{\text{ نماد}}{مفهوم}
\\
%مقادیر بالا را تغییر ندهید
%%%%%%%%%%%%%%%%%%%%%%%%%%%%%%%%%%%%%%%%%%%%%%%%%%%%%%%%%
\symb{\mathbb{R}^n}{
فضای اقلیدسی با بعد $n$
}
\symb{\mathbb{S}^n}{
کره یکه $n$ بعدی
}
\symb{M^m}{
خمینه $m$-بعدی $M$
}
\symb{\mathfrak{X}(M)}{
جبر میدان‌های  برداری هموار روی $M$
}
\symb{\mathfrak{X}^1(M)}{
مجموعه میدان‌های برداری هموار یکه روی $(M,g)$ 
}
\symb{\Omega^p(M)}{
مجموعه $p$-فرمی‌های روی خمینه $M$
}
\symb{Q}{
اپراتور ریچی
}
\symb{\mathcal{R}}{
تانسور انحنای ریمان
}
\symb{ric}{
تانسور ریچی
}
\symb{L}{
مشتق لی
}
\symb{\Phi}{
2-فرم اساسی خمینه تماسی
}
\symb{\nabla}{
التصاق لوی-چویتای
}
\symb{\Delta}{
لاپلاسین ناهموار
}
\symb{\nabla^*}{
عملگر خودالحاق صوری القا شده از التصاق لوی-چویتای
}
\symb{g_s}{
متر ساساکی
}
\symb{\nabla}{
التصاق لوی-چویتای وابسته به متر ساساکی
}
\symb{\Delta}{
عملگر لاپلاس-بلترامی روی $p$-فرم‌ها
}

%%%%%%%%%%%%%%%%%%%%%%%%%%%%%%%%%%%%%%%

\thispagestyle{style3}
\newpage
%\pagestyle{style1}
%%%%%%%%%%%%%%%%%%%%%%%%%%%%%%%%%%%%


\pagenumbering{arabic}
\pagestyle{style1}
%--------------------------------------------------------------------------chapters(فصل ها)
\chapter{مقدمه}

مجموعه داده شامل بیش از 284807 داده از داده‌‌های بانک‌های اروپا شامل بیش از ۲۵ کلاس و همچنین زمان و مقدار آن‌هاست که در اینجا با 1 و 0 شماره گذاری شده اند.

ما در اینجا قصد داریم ابتدا با پیش پردازش و سپس با استفاده از الگوریتم \lr{kmeans} برای خوشه بندی و همچنین الگوریتم های \lr{SVC, LogisticRegression, KNeighborsClassifier} 
به پاسخی برای دسته بندی و همچنین خوشه بندی مجموعه داده برسیم.	

\chapter{}
\section{خوشه بندی و طبقه بندی}
\subsection{خوشه بندی}

خوشه‌بندی یا آنالیز خوشه(\lr{Clustering})در آمار و یادگیری ماشینی، یکی از شاخه های یادگیری بدون ناظر می‌باشد که آن ورودی هست و خروجی ای وجود ندارد و مدل خودش الگوی نهفته داده را پیدا کرده و سپس نمونه‌ها را به دسته‌هایی که اعضای آن مشابه یکدیگر می‌باشند تقسیم می‌کند که به این دسته ها خوشه گفته میشود. بنابراین خوشه مجموعه ای از اشیاء می‌باشد که در آن اشیاء با یکدیگر مشابه بوده و با اشیاء موجود در خوشه‌های دیگر غیر مشابه می‌باشند.

\subsection{طبقه بندی}

طبقه‌ بندی (\lr{classification}) علمی است که بر اساس داده‌های قبلی که دارای برچسب هستند، مدلی برای پیش بینی برچسب داده‌های جدید می‌سازد.


\chapter{
	پیاده سازی
}

\section{مقدمه}

کد به زبان \lr{Python} نوشته شده است، بخش‌های کامل کد رو می‌توانید در  
\href{https://github.com/MohammadrezaAmani/AI_HW5}{گیتهاب}
مشاهده کنید.\\
کد به شکل کامل با استفاده از مفاهیم \lr{OOP} نوشته شده است. لیست کلاس ها شامل موارد زیر می‌باشد:\\


\section{پیش‌نیازها}

برای استفاده از کد ها نیاز به استفاده از زبان برنامه نویسی پایتون و کتابخانه‌های \lr{numpy, pandas, matplotlib, seaborn} و \lr{scikit-learn} دارید.
که با دستور 
\lr{pip install numpy pandas matplotlib seaborn scikit-learn}
میتوانید آن‌ها را نصب کنید.

\section{ساختار کلی کد}
کد به شکل OOP نوشته شده اما در آن متاسفانه زیبایی برنامه نویسی رعایت نشده که ان شاالله در آینده تغییر خواهد یافت.
کد از یک کلاس به نام \lr{CreditCardFraudClassifier} تشکیل شده است که وظیفه ی آن ایمپلیمنت کردن مباحث مورد نیاز برای کد است.

\lr{توابع}

\subsection{\lr{info}}
این تابع اطلاعات کلی از کد را نمایش می‌دهد.

\subsection{\lr{describe}}
مولفه های کد مانند کمینه مقدار و بیشینه مقدار و فیچر ها را نمایش می‌هد.
\subsection{\lr{process}}
کار های پیش پردازش دیتا بر عهده ی این تابع می‌باشد که شامل 
\\ 
بالانس کردن دیتا
\\ 
نورمال کردن مقادیر آن
\\ حذف مقادیر بی تاثیر مانند زمان
\\
جدا کردن داده به نسبت بیست به هشتاد برای رسیدن به پاسخ نهایی و تست آن 

\subsection{\lr{visualize\_kmeans\_clusters}}
این تابع با استفاده از روش خوشه بندی \lr{kmeans}
مجموعه داده را در دو مجموعه خوشه بندی کرده و نتیجه ی نهایی را نمایش می‌دهد.

\subsection{\lr{test\_all\_models}}

این تابع با دریافت لیستی از مدل ها تمام آن‌ها را تست کرده و نتایح را تحلیل و گزارش می‌کند.


\chapter{نتیجه گیری}

\section{تاثیر پیش پردازش}

در روند توسعه ی نرم افزار به این نتیجه رسیدیم که در صورت باقی ماندن دیتا های نادرست باعث به خطا افتادن مدل هوش مصوعی خود میشویم و اگر دیتا را بالانس نکنیم 
مدل خوشه بندی به سمت مدلی که دیتای بیشتری دارد میل می‌کند.

\section{سرعت مدل}
در این کد متوجه می‌شویم مدل رندوم فارست برای این دیتاست بهترین دقت را دارد ولی به دلیل کند بودن آن نمیتوان به آن تکیه کرد.
لیست باقی مدل ها و سرعت و زمان و دقت آنها:

\lr{
RandomForest: 0.99964            3:12
LogisticRegression: 0.99912      0:10
KNeighborsClassifier: 0.99953   0:5
SVC 0.99932                    2:45
}

\chapterfont{\vspace*{-2em}\centering\LARGE}%

\appendix
\bibliographystyle{unsrt-fa}
\bibliography{references}
\end{document}